\section{Musical Primitives and Compositional Constructs}\label{sec:musical-primitives}

To bridge the gap between abstract musical theory and formal computer science, the Motivo language maps
fundamental musical
entities to specific programming constructs.
This mapping ensures that the language remains expressive for musicians while retaining the type safety and structural
clarity of a compiled language.

\subsection{Musical Entities as Data Types}\label{detail-subsec:musical-entities-as-data-types}

The language defines core primitives that serve as the building blocks of a composition.
Each maps to a member of \texttt{types::Type} during semantic analysis.

\noindent \textbf{Notes.} Represented as pitch-class and octave pairs (e.g., \texttt{C4}, \texttt{F\#2}).
These are treated as constant literals with unique numeric identifiers in the underlying MIDI protocol.
Variables may be declared with the \texttt{note} type keyword.

\noindent \textbf{Chords.} Modeled as parenthesized collections of simultaneous notes with optional per-note
durations (e.g.\ \texttt{(C4, E4, G4)}).
A chord is a single playable unit whose temporal span is determined by its constituent notes.

\noindent \textbf{Sequences.} Represented as bracketed ordered lists of notes, chords, or rests with optional
durations (e.g.\ \texttt{[C4, D4:2, rest]}).
This structure maps directly to the concept of a musical ``bar'' or ``cell'' and may be stored in \texttt{seq}
variables or returned from patterns.

\noindent \textbf{Rests.} Unlike many languages that use a \texttt{sleep} function to pause execution, Motivo treats a
rest as a durational musical value (\texttt{rest} or \texttt{rest:duration}).
A rest occupies space on the timeline; it is not a declarable variable type but is valid inside sequences and
\texttt{play} targets.

\noindent \textbf{Drum hits.} Named literals such as \texttt{kick} and \texttt{snare} map to General MIDI percussion
pitches.
They are typed as \texttt{Drum} and are intended for tracks declared with \texttt{using drums}; the semantic pass
rejects melodic notes on drum tracks and drum hits elsewhere.

\subsection{Compositional Abstractions}\label{detail-subsec:compositional-abstractions}

The hierarchical organization of music is reflected in the language's scoping and container constructs:

\subsubsection{Patterns: The Unit of Reusability}
A \texttt{pattern} is the primary abstraction mechanism in Motivo\@.
Semantically, it is equivalent to a function definition.
It allows a composer to encapsulate a motif and assign it a name.
Patterns accept a typed parameter list (e.g.\ \texttt{pattern fill(bool accent, seq motif)}), enabling dynamic musical
ideas---such as a drum fill gated by a \texttt{bool} flag or a melodic line parameterized by a \texttt{seq} template.
Version~2.0 resolves calls by the full type signature, so two patterns named \texttt{transpose} with
\texttt{(int steps)} and \texttt{(double ratio)} can coexist without ambiguity.

\subsubsection{Tracks: Instrumental Identity}
A \texttt{track} serves as the top-level execution context for a specific instrument.
In the compilation process, each track maps to a unique MIDI channel.
This means that independent musical lines (e.g., a piano melody and a bassline) are kept separate and assigned to
distinct sound patches in the final artifact.

\subsubsection{Voices: Horizontal Parallelism}
Within a single track, the \texttt{voice} construct allows for polyphonic composition.
While a track represents a single ``performer'', a voice represents a single ``line of counterpoint''.
Multiple voices can be defined to start at the same time or at different offsets within the same track, allowing for
complex harmonic structures without the need for manual delta-time calculations.
Voices enable the composer to think in terms of horizontal layers rather than just vertical chords.
